\chapter{One Problem, Four Answers}
\label{ch:oneproblem}

\section{The problem}

Fix a symmetric positive definite matrix $G \in \R^{n\times n}$, a vector
$a \in \R^n$, a matrix $C \in \R^{n \times m}$ whose columns $c_1,\dots,c_m$ are
constraint normals, and a right-hand side $b \in \R^m$. The first
$m_{\mathrm{eq}}$ constraints are equalities and the rest inequalities. The
problem is
\begin{equation}
  \min_{x \in \R^n} \; f(x) := \tfrac12 x\T G x - a\T x
  \qquad \text{subject to} \qquad
  \begin{cases}
    c_i\T x = b_i, & i \le m_{\mathrm{eq}},\\
    c_i\T x \ge b_i, & i > m_{\mathrm{eq}},
  \end{cases}
  \label{eq:qp}
\end{equation}
which we will call \emph{the} quadratic program, or QP, throughout.
\index{quadratic program!statement}

Two conventions deserve a word, because they differ from some textbooks and are
inherited from the solver this book's Part~II describes. The linear term is
\emph{subtracted}, so that the unconstrained minimiser is $G^{-1}a$ rather than
$-G^{-1}a$; and the constraints are held \emph{column-wise} in $C$, so the
constraint system reads $C\T x \ge b$ rather than $Cx \le b$. Neither choice
affects anything beyond signs, and both make the formulas in Part~II shorter.

Everything in this book is downstream of two hypotheses in~\eqref{eq:qp}. The
first is that $G \succ 0$, which makes $f$ strictly convex and coercive, so the
problem has exactly one solution whenever it has any. The second is that the
constraints are \emph{linear}, so the feasible set is a polyhedron: an
intersection of finitely many half-spaces and hyperplanes, hence a set with
finitely many faces. Together these give the problem a character that is worth
stating before any algorithm appears.

\begin{quote}
\emph{The analysis is easy and the combinatorics is hard.}
\end{quote}

If an oracle told you which of the $m$ constraints hold with equality at the
solution, you could finish in one linear solve; Chapter~\ref{ch:certificate}
proves this. There are $2^{m-m_{\mathrm{eq}}}$ candidate answers the oracle might
give, and the whole of algorithmic practice is a strategy for searching among
them without enumerating them.

\section{Where it comes from}

The QP~\eqref{eq:qp} is not a mathematical curiosity to be motivated after the
fact. It is solved a great many times a day, in at least four traditions that
have historically not read one another's papers.

\paragraph{Portfolio selection.} An investor spreads weights $w \in \R^n$ across
$n$ assets with covariance $\Sig$ and expected returns $\mu$, subject to being
fully invested, $\onevec\T w = 1$, and---under a long-only mandate---holding no
short positions, $w \ge 0$. Minimising variance against a return tilt $\rho$,
\[
  \min_w \; w\T \Sig w - \rho\,\mu\T w
  \qquad\text{subject to}\qquad \onevec\T w = 1, \quad w \ge 0,
\]
is~\eqref{eq:qp} with $G = 2\Sig$, $a = \rho\mu$, and $C = [\onevec, I]$. This is
Markowitz's problem \cite{markowitz1952}, and it is the running application of
Parts~III and~IV.
\index{portfolio selection}

\paragraph{Non-negative least squares.} Given $M \in \R^{m\times n}$ and
$d \in \R^m$, minimise $\norm{Mx - d}_2^2$ over $x \ge 0$. Expanding the square
gives~\eqref{eq:qp} with $G = 2M\T M$, $a = 2M\T d$, $C = I$, $b = 0$. This is
the problem of Lawson and Hanson \cite{lawson1974}, and it is the setting in
which Chapter~\ref{ch:krylov} works, because $G$ arrives as a factored
Gram matrix and the whole method can be run without ever assembling it.
\index{non-negative least squares}

\paragraph{Model predictive control.} At each sampling instant a controller
minimises a quadratic cost over a finite horizon of inputs subject to linear
state and input constraints. The current state enters as a \emph{parameter} of
the QP. Solving the family once, as a function of that parameter, rather than
solving one instance per instant, is explicit model predictive control
\cite{bemporad2002}, and it is the multi-parameter cousin of the theory in
Part~III.
\index{model predictive control}

\paragraph{Regularised regression.} The LASSO
\cite{tibshirani1996}, $\min_\beta \tfrac12\norm{y - X\beta}_2^2 +
\lambda\norm{\beta}_1$, does not look like~\eqref{eq:qp}: its penalty is not
smooth. But its subgradient condition is a sign pattern on a set of coordinates,
which is an active set by another name, and under the substitution
$G = X\T X$, $a = X\T y$ its solution path is the \emph{same} piecewise-linear
curve as a portfolio frontier. Theorem~\ref{thm:identity} makes this exact. It is
the sharpest illustration in the book of the claim that these are one subject.
\index{LASSO}

The four traditions arrived at the same algorithms independently, four decades
apart in the two extreme cases. Appendix~\ref{app:history} tells that story;
Appendix~\ref{app:notation} translates the vocabulary.

\section{The active set}

\begin{definition}[Active set, working set]
\label{def:activeset}
\index{active set!definition}
For a feasible $x$, the \emph{active set} at $x$ is
$\Active(x) = \{i : c_i\T x = b_i\}$, the constraints holding with equality
there. A \emph{working set} is any index set $\Active \subseteq \{1,\dots,m\}$
that an algorithm is currently \emph{holding} as equalities, whether or not it
is the active set at the current iterate or at the solution.
\end{definition}

The distinction between the two matters and is a frequent source of confusion.
The active set is a property of a point; the working set is a piece of algorithm
state. At a solution the working set of a correct algorithm is a subset of the
active set, and the two coincide except in the degenerate case of
Chapter~\ref{ch:geometry}.

Given a working set $\Active$ of size $k$ with normals $A = C_{\cdot\Active}$, the
\emph{working-set subproblem} drops every constraint outside $\Active$ and
promotes every constraint inside it to an equality:
\begin{equation}
  \min_{x} \; \tfrac12 x\T G x - a\T x
  \qquad \text{subject to} \qquad A\T x = b_\Active .
  \label{eq:sub}
\end{equation}
This is an equality-constrained problem, and equality-constrained problems are
linear algebra. Chapter~\ref{ch:certificate} solves~\eqref{eq:sub} in closed
form, Chapter~\ref{ch:freeblock} shows how to solve it fast and how to update the
solution when $\Active$ changes by one index, and the rest of the book is about
choosing $\Active$.

\section{Four answers}

Here are the four strategies, stated in one paragraph each so that the reader
knows where the book is going. Each gets a chapter, or three.

\paragraph{Walk to it.} Start from a working set you can solve trivially, and
move to the solution one index at a time, maintaining part of the optimality
conditions exactly and driving the violated part to zero. The classical primal
method maintains feasibility and works toward stationarity, which requires a
feasible starting point and therefore a phase-one problem. The \emph{dual}
method of Goldfarb and Idnani \cite{goldfarb1982,goldfarb1983} reverses the
roles: it maintains stationarity and the multiplier signs, and works toward
feasibility. Its starting point is free---with no constraint active, $x = G^{-1}a$
satisfies everything the method asks---so one Cholesky factorisation replaces the
phase-one problem entirely. Chapter~\ref{ch:walking} derives it, and shows that
the whole method follows from two matrix identities.

\paragraph{Guess it.} Solve the subproblem on a guessed working set, read the
signs of the resulting multipliers and slacks, and let those signs tell you which
indices to add and drop---\emph{all of them at once}. This is the primal--dual
active set strategy \cite{hintermuller2002}, equivalently block principal
pivoting on a linear complementarity problem \cite{judice1994,murty1988}. It
converges in a handful of repairs almost independently of $n$, trading arithmetic
(abundant) for iterations (not). Chapter~\ref{ch:guessing} develops it, proves
finite termination in the bound-constrained case $C = I$, and shows exactly why
that theorem does not survive the general constraint class.

\paragraph{Trace it.} Attach a scalar parameter to the linear term,
$a \mapsto \lambda d$, and ask how the solution moves as $\lambda$ varies. The
answer is that it is \emph{piecewise linear}: on any interval where the active set
is constant the optimiser is affine in $\lambda$, and the active set changes at
finitely many breakpoints. So the whole family can be traced exactly, corner to
corner, by a single finite algorithm. Part~III develops this, and shows that
Markowitz's Critical Line Algorithm \cite{markowitz1956}, least angle regression
\cite{efron2004}, the support-vector-machine path \cite{hastie2004svm}, and
explicit model predictive control \cite{bemporad2002} are one algorithm with four
fillings.

\paragraph{Refuse to form the matrix.} Notice how little of $G$ any of the above
actually touches. An active-set method needs a product $v \mapsto Gv$, a solve
against the free block, and a cross product between free and blocked
coordinates---three linear \emph{actions}, not $n^2$ stored numbers. Anything
that supplies those actions can play the role of $G$: a factored Gram matrix
$M\T M$ reached through two products with $M$, a diagonal-plus-low-rank factor
model inverted by the Woodbury identity, a kernel given implicitly by its kernel
function. Chapters~\ref{ch:krylov} and~\ref{ch:operator} develop this reading,
which is where most of the practical speed in this book comes from.

\section{Certificates, and why they are the point}

Two of the four strategies come with finite-termination theorems and two do not.
Guessing can cycle. A parametric trace can be derailed by degeneracy. It would be
natural to conclude that the guaranteed methods are the serious ones and the rest
are heuristics.

That conclusion is wrong, and understanding why is the main methodological lesson
of this book. Because $G \succ 0$, the Karush--Kuhn--Tucker conditions
for~\eqref{eq:qp} are not merely necessary but \emph{sufficient}
(Proposition~\ref{prop:sufficient}). A point equipped with multipliers passing
four inequalities \emph{is} the answer, whatever produced it. So a method that
produces candidates may be as unprincipled as one likes, provided each candidate
is tested; and the test is cheap, being one matrix--vector product and a pass over
the constraints.

This turns the design question around. Instead of asking ``does this method
converge?'' one asks ``can this method recognise its own success, and what does it
do when it fails?'' A method with a certificate and a fallback is safe without
being guaranteed. A guaranteed method whose guarantee rests on an assumption you
cannot check---non-degeneracy, say---may be less safe in practice than an
unguaranteed one that verifies its answer. Chapter~\ref{ch:certificate} proves the
sufficiency; Chapters~\ref{ch:guessing}, \ref{ch:krylov}, and~\ref{ch:warmstart}
each cash it in.

\section{What this book does not cover}

Honesty about scope. The book treats convex quadratic objectives over polyhedra,
and stays there.

\emph{Interior-point methods} are barely mentioned, and only for comparison. They
are the right tool for large sparse problems and they are excellently covered
elsewhere \cite{nocedal2006,boyd2004}. Their relationship to this book's methods is
worth stating once: an interior-point method approaches a face of the polyhedron
asymptotically, whereas an active-set method arrives at one exactly and in finitely
many steps. That is why active-set methods warm-start almost perfectly and
interior-point methods do not, and warm starts are what Part~III lives on.

\emph{Non-convex quadratic programs} ($G$ indefinite) are a different subject:
NP-hard, with local minima, and none of the sufficiency arguments here apply.
\emph{Semidefinite $G$} is a middle case and is genuinely within reach---Boland
\cite{boland1997} extends the dual method to it with finite termination, and the
solver of \cite{arnstrom2022daqp} handles it by proximal regularisation---but the
factorisation of Chapter~\ref{ch:walking} needs $G^{-1}$, so the book keeps
$G \succ 0$ as a matter of scope. Where that restriction bites, the text says so.

\emph{General nonlinear constraints} are out of scope, though the QP is not: a
sequential quadratic programming method solves one instance of~\eqref{eq:qp} per
outer iteration, so everything here is a subroutine of that larger subject.

\begin{notes}
The QP~\eqref{eq:qp} in this form, with its two conventions, is the signature of
the \texttt{cvx-quadprog} package \cite{cvxquadprog}, whose mathematics
Chapters~\ref{ch:walking} and~\ref{ch:guessing} follow. Nocedal and Wright
\cite[Ch.~16]{nocedal2006} is the standard reference for quadratic programming
generally and should be on the desk alongside this book; Boyd and Vandenberghe
\cite{boyd2004} for the convex-analytic background. Golub and Van Loan
\cite{golub2013} is the linear-algebra reference used throughout.

The claim that four fields developed one algorithm independently is argued in
full in \cite{schmelzer2026perspective} and retold in Appendix~\ref{app:history}.
\end{notes}

\begin{exercises}

\begin{exercise}
Write the non-negative least squares problem $\min_{x\ge0}\norm{Mx-d}_2^2$ in the
form~\eqref{eq:qp}, giving $G$, $a$, $C$, $b$, and $m_{\mathrm{eq}}$ explicitly.
For which $M$ is the resulting $G$ positive definite rather than merely
semidefinite?
\end{exercise}

\begin{exercise}
Show that the long-only fully invested portfolio problem has a non-empty feasible
set for every $n \ge 1$, and that the feasible set is compact. Deduce that the
problem has a solution without appealing to coercivity of $f$.
\end{exercise}

\begin{exercise}
The feasible set of~\eqref{eq:qp} with $m_{\mathrm{eq}} = 0$ and $C = I$, $b = 0$
is the non-negative orthant. How many faces does it have? How many candidate
active sets must an enumerating algorithm consider for $n = 50$? Compare with the
number of atoms in the observable universe ($\approx 10^{80}$).
\end{exercise}

\begin{exercise}
Let $x_u = G^{-1}a$ be the unconstrained minimiser. Show that if $x_u$ is
feasible then it solves~\eqref{eq:qp}, and that the active set at the solution is
then contained in $\{i : c_i\T x_u = b_i\}$. Why does this observation make a good
starting heuristic for the guessing strategy of Chapter~\ref{ch:guessing}?
\end{exercise}

\begin{exercise}
Consider~\eqref{eq:qp} with $n = 2$, $G = I$, $a = (0,0)\T$, and the single
constraint $x_1 + x_2 \ge 1$. Solve it by hand, identify the active set at the
solution, and verify that the solution lies on the constraint boundary. Now change
$a$ to $(2,2)\T$ and repeat. What happened to the active set, and why?
\end{exercise}

\begin{exercise}[Implementation]
Implement a brute-force QP solver: enumerate every subset of the inequality
constraints, solve the equality-constrained subproblem~\eqref{eq:sub} on each,
discard the infeasible ones, and return the feasible point of least objective
value. Test it on random instances with $n = 5$ and $m = 8$. Time it as $m$ grows
to $16$. This solver is correct and useless, and the rest of the book explains
why both halves of that sentence are true.
\end{exercise}

\end{exercises}
